\section{Models}
\label{sec:models}

\subsection{Baseline: Long Short-Term Memory (LSTM) Network}

The first model is a Long Short-Term Memory (LSTM) network \parencite{hochreiter1997lstm}, which serves as a deep learning baseline for capturing temporal dependencies in financial return sequences. The architecture contains three stacked LSTM layers with a hidden size of 128 units and a dropout rate of 0.247 applied between layers, both selected via hyperparameter tuning (See table \ref{tab:lstm_tuning}).

The model is trained using Mean Squared Error (MSE) loss. Since standard LSTMs do not inherently provide uncertainty estimates, Monte Carlo (MC) Dropout \parencite{gal2016dropout} is applied at inference time: dropout layers are kept active during evaluation, and $N=50$ stochastic forward passes are performed for each input. The predictive mean and variance are then estimated as:
\begin{equation}
    \hat{\mu} = \frac{1}{N}\sum_{i=1}^{N}\hat{y}^{(i)}, \quad
    \hat{\sigma}^2 = \frac{1}{N}\sum_{i=1}^{N}(\hat{y}^{(i)} - \hat{\mu})^2.
    \label{eq:mc_dropout}
\end{equation}
The LSTM thus acts as a deterministic forecasting model during training and an approximate Bayesian model during inference, with its uncertainty derived from equation~\ref{eq:mc_dropout}.

\subsection{Variational GP Model}

The second model is a Variational Gaussian Process (GP) \parencite{10.7551/mitpress/3206.001.0001, titsias2009variational}, implemented using GPyTorch \parencite{gardner2021gpytorchblackboxmatrixmatrixgaussian}. It provides a fully probabilistic framework for regression with explicit uncertainty quantification. An RBF kernel is used for its ability to capture smooth temporal variations while maintaining computational stability.

To scale to approximately 7,500 training samples, the model uses 100 learnable inducing points \parencite{titsias2009variational} (See table \ref{tab:vgp_tuning}). These act as a compressed representation of the full dataset, enabling efficient approximate posterior inference without requiring full $\mathcal{O}(n^3)$ computation.

Training maximises the ELBO, which jointly optimises data fit and regularisation through a KL divergence term. Unlike the LSTM, the VGP directly outputs a predictive posterior distribution:

\begin{equation}
    p(y_* \mid x_*, X, y) = \mathcal{N}(\mu_*, \sigma_*^2),
    \label{eq:vgp_posterior}
\end{equation}

where both mean and variance in equation \ref{eq:vgp_posterior} are obtained analytically from the approximate GP posterior.

\subsection{Metrics}
\label{sec:metrics}

Model performance is evaluated using both predictive accuracy and uncertainty quality metrics. Root Mean Squared Error (RMSE) measures point prediction accuracy as equation \ref{eq:rmse}:
\begin{equation}
    \mathrm{RMSE} = \sqrt{\frac{1}{N}\sum_{i=1}^N (y_i - \hat{y}_i)^2}.
    \label{eq:rmse}
\end{equation}

Hit Rate evaluates directional accuracy by measuring how often the model correctly predicts the sign of the return as equation \ref{eq:hit_rate}:
\begin{equation}
    \mathrm{HitRate} = \frac{1}{N}\sum_{i=1}^N \mathbf{1}\!\left[\mathrm{sign}(\hat{y}_i) = \mathrm{sign}(y_i)\right].
    \label{eq:hit_rate}
\end{equation}

Since both models output Gaussian predictive distributions $\mathcal{N}(\mu_i, \sigma_i^2)$, the Negative Log-Likelihood (NLL) is equation \ref{eq:nll}:
\begin{equation}
    \mathrm{NLL} = \frac{1}{N}\sum_{i=1}^N \left[\frac{(y_i - \mu_i)^2}{2\sigma_i^2} + \frac{1}{2}\log(2\pi\sigma_i^2)\right].
    \label{eq:nll}
\end{equation}
NLL is a proper scoring rule that simultaneously penalises inaccurate means and miscalibrated variances.

Finally, the Expected Calibration Error (ECE) quantifies miscalibration by comparing predicted confidence with empirical coverage across $B$ equal-width bins, as equation~\ref{eq:ece}:
\begin{equation}
    \mathrm{ECE} = \sum_{b=1}^B \frac{|B_b|}{N} \left|\mathrm{acc}(B_b) - \mathrm{conf}(B_b)\right|,
    \label{eq:ece}
\end{equation}
where $\mathrm{acc}(B_b)$ is the empirical coverage in bin $b$ and $\mathrm{conf}(B_b)$ is the mean predicted confidence.


\subsection{Hyperparameter Tuning and Training Configuration}

Hyperparameters were selected using Optuna \parencite{optuna_2019} with a GP sampler, minimising validation NLL over 50 trials (10 complete, 40 pruned via median pruning) across 5-fold cross-validation, with a training budget of 50 epochs per trial (patience 10). Full search spaces and best values are reported in appendix \ref{app:hyperparameter_tuning}.

For the LSTM, the search favoured a deeper architecture (3 layers, hidden size 128) with moderate dropout (0.247) and a relatively small batch size (32). For the VGP, a higher learning rate ($1.45 \times 10^{-2}$) and 100 inducing points were selected. Cross-validation NLL values are not directly comparable between models: the LSTM is trained with MSE loss and its NLL is evaluated post-hoc, while the VGP optimises the ELBO directly.

Final models were trained for 100 epochs with early stopping (patience 10), a fixed random seed of 33, and a CUDA-enabled GPU. The best checkpoint (by validation NLL) was restored after early stopping.

